\chapter{Auswertung}
\label{ch:Auswertung}
% Liste der genutzer Formeln für die Fehlerrechnung
\section*{Fehlerrechnung}
Für die statistische Auswertung von $n$ Messwerten $x_i$ werden folgende Größen definiert \cite{errorSkript25}:
\begin{align}
    \bar{x} &= \frac{1}{n} \sum_{i=1}^{n} x_i \vphantom{\sqrt{\sum_i^n}^2} && \text{\textcolor{gray}{Arithmetisches Mittel}} \label{eq:arithmetisches_mittel} \\
    \sigma^2 &= \frac{1}{n-1} \sum_{i=1}^{n} (x_i - \bar{x})^2 \vphantom{\sqrt{\sum_i^n}^2} && \text{\textcolor{gray}{Varianz}} \label{eq:Varianz} \\
    \sigma &= \sqrt{\frac{1}{n-1} \sum_{i=1}^{n} (x_i - \bar{x})^2} \vphantom{\sqrt{\sum_i^n}^2} && \text{\textcolor{gray}{Standardabweichung}} \label{eq:standardabweichung} \\
    \Delta \bar{x} &= \frac{\sigma}{\sqrt{n}} = \sqrt{\frac{1}{n(n-1)} \sum_{i=1}^n(\bar x - x_i)^2} \vphantom{\sqrt{\sum_i^n}^2} && \text{\textcolor{gray}{Fehler des Mittelwerts}} \label{eq:fehler_mittelwert} \\
    \Delta f &= \sqrt{\left(\frac{\partial f}{\partial x} \Delta x\right)^2 + \left(\frac{\partial f}{\partial y} \Delta y\right)^2} \vphantom{\sqrt{\sum_i^n}^2} && \text{\textcolor{gray}{Gauß’sches Fehlerfortpflanzungsgesetz für $f(x,y)$}} \label{eq:gauss_fehlfortpflanzung} \\
    \Delta f &= \sqrt{(\Delta x)^2 + (\Delta y)^2} \vphantom{\sqrt{\sum_i^n}^2} && \text{\textcolor{gray}{Fehler für $f = x + y$}} \label{eq:fehler_summe} \\
    \Delta f &= |a| \Delta x \vphantom{\sqrt{\sum_i^n}^2} && \text{\textcolor{gray}{Fehler für $f = ax$}} \label{eq:fehler_proportional} \\
    \frac{\Delta f}{|f|} &= \sqrt{\left(\frac{\Delta x}{x}\right)^2 + \left(\frac{\Delta y}{y}\right)^2} \vphantom{\sqrt{\sum_i^n}^2} && \text{\textcolor{gray}{relativer Fehler für $f = xy$ oder $f = x/y$}} \label{eq:relativer_fehler} \\
    \sigma &= \frac{|a_{lit} - a_{gem}|}{\sqrt{\Delta a_{lit}^2 + \Delta a_{gem}^2}} \vphantom{\sqrt{\sum_i^n}^2} && \text{\textcolor{gray}{Berechnung der signifikanten Abweichung}} \label{eq:signifikante_abweichung}
\end{align}

\newpage

\section{Bestimmung der Halbwertszeit und der Dämpfung}
\label{sec:A1}
Eine \afz{optimale} Kreiselbewegung ist nicht möglich, es treten Reibungskräfte auf, die die Rotation beeinflussen. Diese wird als Dämpfung bezeichnet. 
Die Dämpfung im System bleibt dabei kontinuirlich die Gleiche, und wird daher hier auch als \emph{Dämpfungskonstante} bezeichnet.
Jene Dämpfung für das vorliegende System kann über die Steigung der Regression nach \cabb{Damping_1} bestimmt werden. Die benutzen Messdaten sind dem \hyperref[ch:Protokoll]{Messprotokoll} zu entnehmen. 

Die Messdaten wurden jedoch nicht 1:1 übernommen, die Frequenzen mussten in Kreiselfrequenzen umgerechtent werden:
\eq{f_to_omega}{
    \omega = 2\pi\, f \, , \quad \Delta\omega = 2\pi \, \Delta f \, .
}

Es ist auch zu beachten, dass die Frequenz nicht linear abfällt, sondern in einem Exponentialzusammenhang steht:
\eq{exp_zusammenhang}{
    \omega (t) = \omega_0 \, \exp{[-\delta \, t]} \, .
}

$\omega_0$ beschreibt die Startfrequenz, $t$ die verstrichene Zeit und $\delta$ die Dämpfung.

\img{Damping_1}{Kreisfrequenz $\omega_F$ gegen die Zeit $t$ aufgetragen. Die Ordinate ist logarithmisch Skaliert.}

Durch die lograrithmische Skalierung erscheint die Regression linear. Für die Steigung der Regressionsfunktion ergibt sich:

\res[-4]{\delta}{7.82}{0.16}{s^{-1}} \label{eq:delta_1}

Über die Dämpfung lässt sich auch die Halbwertszeit $T_\text{H}$ bestimmen. Die logarithmische Skalierung msus beachtet werden.

\eq{halb_zeit_f}{
    T_\text{H} = \frac{\ln{2}}{\delta} \, , \quad \Delta T_\text{H} = \frac{\Delta \delta}{\delta}
}

Die Ergebnisse von \ceq{delta_1} werden in \ceq{halb_zeit_f} eingesetzt und berechnet:

\res[2]{T_\text{H}}{8.86}{0.18}{s}

Dies entspricht einer Halbwertszeit von $T_\text{H} = (14.8 \pm 0.3)$ min. 

\section{Vergleich mit Präzessionszeiten}
\label{sec:A2}
Für das weitere Fortgehen muss überprüft werden, wie sich der Auslenkungswinkel $\phi$ auf die Präzessionsdauer $T_P$ auswirkt.

\begin{table}[h]
    \centering
    \caption{Gemessene Präzessionszeiten bei angepassten Winkeln des Kreisels zur Achse}
    \label{tab:P_time}
    \begin{tabular}{L  L}
        \toprule 
        \phi [\mathrm{^\circ}] & T_P [\mathrm{s}] \\
        \midrule 
        45 & 70,3 \pm 0,3 \\ 
        90 & 69,6 \pm 0,3 \\ 
        \bottomrule 
    \end{tabular}
\end{table}

Die Abweichung der beiden Werte liegt bei $\sigma = 1,6$ und ist damit nicht signifikant. 

\section{Berechnung des Trägheitsmomentes in z-Richtung}
\label{sec:A3}

\section{Qualitativer Vergleich der Beobachtung mit der Erwartung}
\label{sec:A4}
% Vergleich von beobachter Nutation mit "Gleichung 8":

% Betrachtung von b = omega_F/Omega: 
% 1. 0 > b >1: I_z < I_x
% 2. Gleich 1: Mathematisch problematisch, im Limes I_x = infty I_z
% 3. Gleich 2: I_z = I_x
% 4. Größer 2: I_z > I_x

% Unsere Beobachtung: Rot in Uhrzeigersinn, aber Farbverlauf entgegen Uhrzeigersinn, daher:

% Omega < omega_f => I_z < I_x (Fall 1)

% Auch erwartbar, denn Figurenachse (die Stange) nur kleiner Teil der Rotation zum I_z zu sich slebst, dafür größer bei I_x = I_y. 


\section{Bestimmung des Trägheitmoments in x-Richtung}
\label{sec:A5}

% Ziel: Aus Drehfrequent der momentaren Drehachse omega_f -> Trägheitsmoment I_x

% Plot: x = omega_f, y = Omega

% Omega = 20pi/t_10
% omega_f wurde gemessen

% $\Delta I_x = I_x \, \sqrt{\frac{\Delta \left(\frac{\Omega}{\omega_F}\right)^{2}}{\left(\frac{\Omega}{\omega_F} - 1\right)^{2}} + \frac{\Delta I_{z}^{2}}{I_{z}^{2}}}$
